\documentclass[12pt,a4paper]{article}
\usepackage[utf8]{vietnam}
\usepackage{amsmath,amssymb}
\usepackage{geometry}
\geometry{
  top=12mm,
  bottom=12mm,
  left=14mm,
  right=14mm,
  footskip=6mm,
  headheight=14pt
}

\usepackage{tikz}
\usetikzlibrary{arrows.meta,calc,shapes.geometric}
\usepackage[most]{tcolorbox}
\usepackage{fontawesome}
\usepackage{fancyhdr}
\usepackage{enumitem}

\pagestyle{fancy}
\fancyhf{}
\fancyhead[L]{\footnotesize\color{teal!80!black}\faBook\ \textbf{SỔ TAY KIẾN THỨC TOÁN 3}}
\fancyhead[R]{\footnotesize\color{teal!80!black}\textbf{Lớp Toán Thầy Lĩnh -- loptoanthaylinh.vn}}
\fancyfoot[C]{\footnotesize\color{gray}-- Sổ tay kiến thức: Bài 21. Khối lập phương, khối hộp chữ nhật --}
\renewcommand{\headrulewidth}{0.4pt}
\renewcommand{\footrulewidth}{0pt}

\setlength{\parindent}{0pt}
\setlength{\parskip}{0pt}

\begin{document}

% --- BANNER BÀI HỌC ---
\begin{tcolorbox}[
  enhanced,
  colback=teal!8!white,
  colframe=teal!70!black,
  arc=3mm, boxrule=1.2pt,
  left=3mm, right=3mm, top=2mm, bottom=2mm,
  halign=center
]
  {\LARGE\bfseries\color{teal!90!black}\faBookmark\ BÀI 21: KHỐI LẬP PHƯƠNG, KHỐI HỘP CHỮ NHẬT}\\[1mm]
  {\itshape\footnotesize\color{black!75} Giải phẫu số đỉnh, số mặt, số cạnh của hai hình khối không gian cơ bản trên 1 trang A4}
\end{tcolorbox}

\vspace{-2mm}

\noindent
\begin{minipage}[t]{0.49\textwidth}\vspace{0pt}
\begin{tcolorbox}[
  enhanced,
  colback=blue!4!white,
  colframe=blue!75!black,
  arc=2mm, boxrule=0.8pt,
  left=2mm, right=2mm, top=1.8mm, bottom=1.8mm,
  title={\footnotesize\bfseries\color{white}\faCube\ 1. Khối hộp chữ nhật}
]
{\footnotesize
\begin{itemize}\setlength{\itemsep}{1pt}\setlength{\leftskip}{-2.5mm}
  \item Gồm có: $\mathbf{8}$ \textbf{đỉnh}, $\mathbf{12}$ \textbf{cạnh} và $\mathbf{6}$ \textbf{mặt}.
  \item Cả $6$ mặt đều là \textbf{hình chữ nhật}.
  \item Các mặt đối diện thì bằng nhau:
  \begin{itemize}\setlength{\itemsep}{0.5pt}\setlength{\leftskip}{-2.5mm}
    \item Mặt đáy trên bằng mặt đáy dưới.
    \item Hai mặt bên trái - phải bằng nhau.
    \item Mặt trước bằng mặt sau.
  \end{itemize}
  \item Ví dụ thực tế: Bao diêm, viên gạch, hộp bánh.
\end{itemize}
}
\end{tcolorbox}
\end{minipage}
\hfill
\begin{minipage}[t]{0.49\textwidth}\vspace{0pt}
\begin{tcolorbox}[
  enhanced,
  colback=teal!4!white,
  colframe=teal!75!black,
  arc=2mm, boxrule=0.8pt,
  left=2mm, right=2mm, top=1.8mm, bottom=1.8mm,
  title={\footnotesize\bfseries\color{white}\faCube\ 2. Khối lập phương (Đặc biệt)}
]
{\footnotesize
\begin{itemize}\setlength{\itemsep}{1pt}\setlength{\leftskip}{-2.5mm}
  \item Gồm có: $\mathbf{8}$ \textbf{đỉnh}, $\mathbf{12}$ \textbf{cạnh} và $\mathbf{6}$ \textbf{mặt}.
  \item Cả $6$ mặt đều là \textbf{hình vuông bằng nhau}.
  \item Toàn bộ $12$ cạnh đều \textbf{dài bằng nhau}.
  \item Khối lập phương là khối hộp chữ nhật đặc biệt có chiều dài, chiều rộng và chiều cao bằng nhau.
  \item Ví dụ thực tế: Con xúc xắc, khối ru-bich.
\end{itemize}
}
\end{tcolorbox}
\end{minipage}

\vspace{-2.5mm}
% --- KHỐI 3: SƠ ĐỒ TƯ DUY TIKZ ---
\noindent
\begin{tcolorbox}[
  enhanced,
  colback=orange!4!white,
  colframe=orange!80!black,
  arc=2mm, boxrule=0.8pt,
  left=2mm, right=2mm, top=1.5mm, bottom=1.5mm,
  title={\footnotesize\bfseries\color{white}\faSitemap\ SƠ ĐỒ TƯ DUY: BẢN ĐỒ CẤU TRÚC 8 ĐỈNH - 12 CẠNH - 6 MẶT}
]
\centering
\begin{tikzpicture}[scale=0.9]
  % Hop chu nhat 3D
  \draw[thick, blue!80!black] (0,0) rectangle (2.8, 1.4);
  \draw[thick, blue!80!black] (2.8,0) -- (3.6, 0.6) -- (3.6, 2.0) -- (2.8, 1.4);
  \draw[thick, blue!80!black] (0, 1.4) -- (0.8, 2.0) -- (3.6, 2.0);
  \draw[dashed, gray] (0,0) -- (0.8, 0.6) -- (3.6, 0.6);
  \draw[dashed, gray] (0.8, 0.6) -- (0.8, 2.0);
  \node[below, font=\scriptsize\bfseries, blue!80!black] at (1.8, -0.2) {Khối hộp chữ nhật (6 mặt CN)};

  % Lap phuong 3D
  \draw[thick, teal!80!black] (7.0,0) rectangle (8.6, 1.6);
  \draw[thick, teal!80!black] (8.6,0) -- (9.4, 0.6) -- (9.4, 2.2) -- (8.6, 1.6);
  \draw[thick, teal!80!black] (7.0, 1.6) -- (7.8, 2.2) -- (9.4, 2.2);
  \draw[dashed, gray] (7.0,0) -- (7.8, 0.6) -- (9.4, 0.6);
  \draw[dashed, gray] (7.8, 0.6) -- (7.8, 2.2);
  \node[below, font=\scriptsize\bfseries, teal!80!black] at (8.2, -0.2) {Khối lập phương (6 mặt vuông)};
\end{tikzpicture}
\end{tcolorbox}

\vspace{-2.5mm}
% --- KHỐI 4: VÍ DỤ & MẸO TRÁNH BẪY ---
\noindent
\begin{minipage}[t]{0.49\textwidth}\vspace{0pt}
\begin{tcolorbox}[
  enhanced,
  colback=green!4!white,
  colframe=green!60!black,
  arc=2mm, boxrule=0.8pt,
  left=2mm, right=2mm, top=1.8mm, bottom=1.8mm,
  title={\footnotesize\bfseries\color{white}\faCheckCircle\ Ví dụ thực hành chuẩn mực}
]
{\footnotesize
\textbf{1. Đếm số khối nhỏ:}\\
Một khối lập phương lớn được ghép từ các khối lập phương cạnh $1\text{ cm}$. Nếu mỗi hàng có $3$ khối, có $3$ hàng và $3$ tầng thì tổng số khối là: $3 \times 3 \times 3 = \mathbf{27}$ (khối).\\
\textbf{2. Nhận biết mặt đối diện:}\\
Trên con xúc xắc, mặt $1$ chấm đối diện với mặt $6$ chấm (tổng số chấm hai mặt đối diện bằng $7$).
}
\end{tcolorbox}
\end{minipage}
\hfill
\begin{minipage}[t]{0.49\textwidth}\vspace{0pt}
\begin{tcolorbox}[
  enhanced,
  colback=red!4!white,
  colframe=red!70!black,
  arc=2mm, boxrule=0.8pt,
  left=2mm, right=2mm, top=1.8mm, bottom=1.8mm,
  title={\footnotesize\bfseries\color{white}\faExclamationTriangle\ Mẹo tránh bẫy trắc nghiệm}
]
{\footnotesize
\begin{itemize}\setlength{\itemsep}{1.5pt}\setlength{\leftskip}{-2.5mm}
  \item \textbf{Bẫy nhầm lẫn số mặt và số đỉnh:}\\
  Hãy nhớ quy tắc bộ 3 số thần thánh: $\mathbf{6}$ \textbf{mặt} -- $\mathbf{8}$ \textbf{đỉnh} -- $\mathbf{12}$ \textbf{cạnh}. Rất nhiều bạn nhớ nhầm thành 8 mặt 6 đỉnh!
  \item \textbf{Cạnh bằng nhau:}\\
  Chỉ có khối lập phương thì $12$ cạnh mới bằng nhau. Khối hộp chữ nhật có $3$ kích thước khác nhau: chiều dài, chiều rộng và chiều cao.
\end{itemize}
}
\end{tcolorbox}
\end{minipage}


\end{document}
